%-*-latex-*-
%-*-latex-*-
\newcommand\COURSE{ciss350}
\newcommand\ASSESSMENT{q03}
\newcommand\ASSESSMENTTYPE{Quiz}
\newcommand\POINTS{\textwhite{xxx/xxx}}

\input{myquizpreamble}
\input{yliow}
\input{\COURSE}
%-*-latex-*-

%-*-latex-*-

%-*-latex-*-
\renewcommand\TITLE{\ASSESSMENTTYPE \ \ASSESSMENT}

\renewcommand\EMAIL{}
\renewcommand\AUTHOR{eagraham1@cougars.ccis.edu}

\textwidth=6in
\begin{document}
\topmatterthree

For both questions, you are writing functions, not a complete programs.

TURN PAGE

%------------------------------------------------------------------------------
\newpage
\nextq
Write a function that performs bubblesort using pointers and recursion
that has the same behavior as our usual bubblesort
(as specified in my 240 notes).
Calling \verb!bubblesort(start, end)! will perform bubblesort (in ascending
order) on the values at address \verb!start! up to \verb!end - 1!.
If \verb!x[0..9]! is an array of size 10, then calling
\verb!bubblesort(&x[0], &x[10])! will sort \verb!x[0..9]!.
You cannot use for-, while-, do-while loops.
You can only use recursion.
A rough skeleton is given.

\ANSWER
\begin{answercode}
void bubblesort_pass(int * start, int * end)
{
    if (start + 1 == end)
    {
        return;
    }
    else
    {
        if (*start > *(start + 1)
        {
          int temp = *start;
          *start = *(start + 1);
          *(start + 1) = temp;
        }
        bubblesort_pass(start + 1, end)
    }
}

void bubblesort(int * start, int * end)
{
    if (start + 1 == end)
    {
        return;
    }
    else
    {
        bubblesort_pass(start, end);
        bubblesort(start, end - 1);
    }
}
\end{answercode}

TURN PAGE

%------------------------------------------------------------------------------
\newpage
\nextq
Write a function that performs binary search (using pointers and recursion)
on the values at address \verb!start! up to \verb!end - 1!.
The values are assumed to be in ascending order.
If \verb!x[0..9]! is an array of size 10, then calling
\verb!binarysearch(&x[0], &x[10], 42)! will return the \textit{address} of 42
in \verb!x[0..9]!.
You cannot use for-, while-, do-while loops.
You can only use recursion.
A rough skeleton is given.

\ANSWER
\begin{answercode}
int * binarysearch(int * start, int * end, int target)
{
    if (start + 1 = end)
    {
        return NULL; // target not found
    }
    else
    {
      int * mid = start + (end - start) / 2;

      if(*mid == target)
        return mid;
      else if(*mid > target)
        return binarysearch(start, mid, target);
      else
        return binarysearch(mid + 1, end, target)
    }
}
\end{answercode}

%------------------------------------------------------------------------------
%-*-latex-*-
\newpage
\input{instructions}
\end{document}

